\chapter{Computing with DNA}\label{ch:computing}
\ChapterGoals{%
\item explain what makes a physical process a \emph{computation}: states, rules and an interpretation;
\item describe how a graph problem can be written into DNA strands, and why strand direction matters;
\item trace a generate--filter--verify search from molecules to a checked answer;
\item distinguish a population of molecules from a set of candidates, and a signal from a proof;
\item estimate how many molecules a search needs, and explain why molecular parallelism does not remove exponential cost.}
{secondary-school biology and algebra; the idea of a graph (Appendix~\ref{app:math} reviews it).}

Can a collection of molecules represent possible solutions to a problem,
undergo operations that remove unsuitable candidates, and leave behind a
readable answer? In 1994 Leonard Adleman showed that the answer is yes for a
small route-finding problem.\cite{Adleman1994} This chapter rebuilds the logic
of that idea from the ground up, connects each logical step to a molecular
mechanism, and then counts the resources the construction demands. Along the
way it introduces the three habits of mind that run through the whole book:
separate the specification from the mechanism and from the evidence; count
copies, not just identities; and never mistake a model's prediction for a
measurement.

\section{Computation in matter}\label{sec:matter}
A computer needs three things: distinguishable \emph{states}, \emph{rules} that
change those states, and an \emph{interpretation} that says what a final state
means. None of these requires the states to be voltages. A DNA sequence can
represent an object; selective pairing between strands can recognise part of
that representation; and a test tube can hold an enormous population of
different strands at once. The interesting step is not writing computer symbols
with four letters. It is making molecular interactions perform an operation
whose result has a precise computational meaning.

Figure~\ref{fig:field} separates three levels that recur throughout this book.
At the \emph{formal} level, a sequence of graph vertices either satisfies a
specification or it does not. At the \emph{physical} level, molecules associate,
dissociate, react and produce measurable signals. At the \emph{interpretive}
level, we decide what those signals establish. Confusing these levels makes a
visually persuasive experiment look more conclusive than it is.

\TikzPlate{chapters/ch01/figures/three-levels}{Three levels of a DNA computation.
The formal level (top) defines what counts as a correct answer for an instance
$X$. The physical level (middle) is an imperfect chemical process acting on a
population of molecules $M$. The interpretive level (bottom) turns an observation
$O$ into a claim, which must be checked against the specification. Read the
arrows as a loop: encode, react, measure, interpret and verify. The drawing is a
conceptual map, not an experimental result.}{fig:field}

Written as a chain,
\[
 X\xrightarrow{\ \mathrm{encode}\ }M
 \xrightarrow{\ \mathrm{physical\ process}\ }O
 \xrightarrow{\ \mathrm{interpret}\ }\text{answer about }X.
\]
Here $X$ is an instance of a problem, $M$ is a prepared molecular population and
$O$ is an observation. An encoding is not yet an implementation of the middle
arrow. Conversely, an impressive chemical transformation is not yet a
computation unless we can say which function of $X$ it computes. The design
problem is to make the whole chain reliable enough for its intended use.

This chapter enters the subject through a graph search. The companion program
models the logical predicates exactly; it does not model reaction rates, which
are the subject of Chapters~6, 7 and~19.

\section{Adleman's experiment: the founding idea}\label{sec:history}
Adleman's 1994 experiment solved a seven-vertex instance of the
\emph{directed Hamiltonian path} problem: is there a route through a network of
one-way roads that starts at a given city, ends at another and visits every
city exactly once? Each vertex was assigned a 20-nucleotide code. Edge
fragments and complementary vertex ``splint'' strands were mixed so that pairing
and an enzyme (ligase) joined fragments into long strands representing random
walks through the graph. The polymerase chain reaction (PCR) then copied only
walks with the right start and end, a gel kept only strands of the length of a
seven-vertex path, and repeated affinity purifications kept only strands that
contained every vertex. Graduated PCR and gel patterns supported the route
$0,1,2,3,4,5,6$.\cite{Adleman1994} This was neither modern sequencing nor a
general efficient algorithm for large graphs.

The experiment made a change of substrate concrete: a combinatorial search could
be organised as preparation, molecular transformation, selection and
interpretation. Its significance does not depend on beating an electronic
computer on seven vertices; it lies in showing that the proposed mapping
survives contact with real chemistry. Chapter~2 reconstructs the historical
experiment in detail.

Here we use an original six-vertex graph instead. It is deliberately different
from the historical instance: it has two Hamiltonian paths, a cycle, short
routes and a useful false candidate. We use it to derive the logic and expose
assumptions. The sequence domains in our figures are symbolic, not validated
laboratory oligonucleotides, and our numerical tables are program outputs, not
measurements.

\section{The graph problem, precisely}\label{sec:graph}
Begin with Figure~\ref{fig:graph}. An arrow from $B$ to $C$ permits that step;
it does not automatically permit a step from $C$ to $B$. In this particular graph
both arrows exist, so a candidate may revisit that pair indefinitely unless
another constraint stops it.
\DeepFigure{DNAD-01-F2}{The original teaching graph. Its vertices are $A$ through
$F$, with fixed start $A$ and finish $F$. The two arrows between $B$ and $C$
create a directed cycle. The graph has nine edges and exactly two Hamiltonian
paths with these endpoints.}{fig:graph}
Formally, $G=(V,E)$, where
\[
 V=\{A,B,C,D,E,F\},\qquad
 E=\{AB,AC,BC,BD,CB,CD,DE,DF,EF\}.
\]
The shorthand $AB$ means the ordered pair $(A,B)$, not an undirected connection.

\begin{definition}
A directed \Term{walk} $P=(v_1,\ldots,v_k)$ has each $v_i\in V$ and each
$(v_i,v_{i+1})\in E$. A simple directed path is a walk with no repeated vertex.
A \Term{Hamiltonian path} visits every vertex exactly once. With fixed endpoints
$s,t$ and $n=|V|$, it satisfies
\[
 k=n,\quad v_1=s,\quad v_n=t,\quad
 v_i\ne v_j\ (i\ne j),\quad
 (v_i,v_{i+1})\in E\ (1\le i<n).
\]
\end{definition}

We keep the word \emph{walk} for candidates that may repeat a vertex. This
matters because a molecule does not know that revisiting $B$ is forbidden merely
because our problem uses the word path.

For $ABCDEF$, each adjacent pair is an edge and all six vertices occur once. For
$ABCBDF$, all adjacent pairs are still edges, the endpoints are right and the
length is six, but $B$ occurs twice and $E$ is absent. For $ABDEF$, the edges
and endpoints are legal but $C$ is missing. For $ABECDF$, all vertices occur
once, yet $BE$ and $EC$ are not edges.
\DeepFigure{DNAD-01-F3}{Four candidates isolate four logical issues: a witness, a
repeated vertex, insufficient length and illegal adjacency. A presence test
cannot detect an illegal edge merely by checking that every vertex
occurs.}{fig:routes}

The \Term{decision problem} asks whether a witness exists. The \emph{search
problem} asks for one. Once a candidate order is available, checking it is
easy; finding that order without being given it is the difficult part. We must
not transfer the ease of verification to the cost of producing a witness.

\section{DNA in one picture}\label{sec:primer}
To write a graph into molecules we need only four facts about DNA, illustrated
in Figure~\ref{fig:primer-helix}. Chapter~5 develops them properly, and
Appendix~\ref{app:biology} gives a fuller primer.

\TikzPlate{chapters/ch01/figures/primer-helix}{DNA as a molecule and as this book
draws it. (A)~A stretch of B-form double helix with approximately correct
proportions: two antiparallel sugar--phosphate backbones (blue and violet) wind
around one axis, and coloured rungs are base pairs (A green, T red, G amber, C
blue). (B)~Unwound, each strand is a covalent chain of sugars (pentagons) and
phosphates (dots); A pairs with T through two hydrogen bonds and G with C
through three. (C)~The abstraction used throughout the book: a thick line with
a half-arrowhead at its $3'$ end, named \emph{domains} for sequence segments,
and a star for the complementary segment. The abstraction keeps orientation and
complementarity; it omits sequence, energy and three-dimensional
shape.}{fig:primer-helix}

\begin{enumerate}
\item \textbf{A strand has a direction.} The sugar--phosphate backbone links the
$3'$ carbon of one sugar to the $5'$ carbon of the next, so the two ends of a
strand are chemically different. Sequences are written from the $5'$ end to
the $3'$ end.
\item \textbf{Complementary strands pair antiparallel.} A pairs with T and G with
C, and the two strands of a duplex run in opposite directions.\cite{WatsonCrick1953,Alberts2002}
\item \textbf{Pairing is reversible; backbone bonds are not.} Heating separates
paired strands and cooling lets them re-pair. Joining two strands into one
requires an enzyme, a \Term{ligase}, to make a new covalent bond.
\item \textbf{A tube holds many copies.} A microlitre of a micromolar solution
contains about $6\times10^{11}$ molecules (Appendix~\ref{app:biology}).
\end{enumerate}

\begin{BioBridge}
Think of a strand as a string with an arrow on it, and of pairing as
two such strings zipping together only if they are complementary \emph{and}
point in opposite directions. The partner of $5'$-ACGA-$3'$, written in its own
$5'\to3'$ direction, is $5'$-TCGT-$3'$: complement each base \emph{and} reverse
the order. This \Term{reverse complement} is the single most common source of
design errors in molecular programming.
\end{BioBridge}

\section{From adjacency to oriented strands}\label{sec:encoding}
In our diagrams, letters such as $b_L$ denote whole sequence domains, not
individual bases. A domain is a contiguous sequence chosen to play a particular
recognition role. Split each vertex code into two equal-length domains:
\[
 O_v=L_vR_v,\qquad |L_v|=|R_v|=h .
\]
An internal edge $u\to v$ is represented by the $5'\to3'$ fragment $R_uL_v$:
the second half of $u$'s code followed by the first half of $v$'s. The next edge
$v\to w$ uses $R_vL_w$. A strand complementary to the vertex code $L_vR_v$ can
pair with the end of the first fragment and the beginning of the second, holding
them side by side (Figure~\ref{fig:encoding}).

\TikzPlate{chapters/ch01/figures/junction}{Domain-level edge encoding at vertex
$B$. (A)~Two edge strands, for $A\to B$ ($a_Rb_L$) and $B\to C$ ($b_Rc_L$), are
written $5'\to3'$. The vertex splint for $B$ runs antiparallel underneath, so
that $b_L^{*}$ pairs with $b_L$ and $b_R^{*}$ with $b_R$. Pairing places the $3'$
end of the first strand beside the $5'$ end of the second, leaving a
\emph{nick}. (B)~Ligase seals the nick, producing one continuous strand.
Domain widths are schematic.}{fig:encoding}

This orientation prevents a common explanatory reversal. In this construction,
complementary \emph{vertex} strands align \emph{edge} fragments; it is not a
picture of whole-vertex strands being glued together by a vague edge label.

Endpoint fragments need a convention too. Complete the first fragment with all
of $O_s$ and the last with all of $O_t$, rather than only their interior-facing
half-domains. For a route with $k$ vertices, the resulting idealised single
strand is $O_{v_1}\cdots O_{v_k}$, of length $2hk$ nucleotides. Length then
becomes a proxy for the number of vertex occurrences; without the endpoint
convention, the length formula changes.

This construction explains the intended correspondence, not sequence-design
success. Unwanted partial matches, folded structures and different binding
strengths can break the clean symbolic picture. Chapters~5, 6 and~11 replace the
domain abstraction with physical constraints.

\section{A population is not an exhaustive search}\label{sec:population}
The candidate pool is best understood as a \Term{multiset}: a route may have many
molecular copies, another may have one, and a valid route may be absent. A
million copies of one wrong route do not cover a million distinct possibilities.
\DeepFigure{DNAD-01-F5}{An illustrative candidate population. Copy count,
distinct-route count and coverage of valid solutions are different quantities.
The last route is an intentionally injected pseudo-path for testing the
verification boundary; the pictured multiplicities are not experimental
abundances.}{fig:pool}

Let $N(P)$ be the number of copies of candidate $P$. The total population is
$\sum_P N(P)$, whereas its \emph{support} is $\{P:N(P)>0\}$. Selection can enrich a
member of the support. In a pure selection model it cannot recover a witness that
was never present.

\begin{Pitfall}
``The tube contains $10^{14}$ strands, so it explores $10^{14}$ candidates.''
Not necessarily. Copies of the same walk add mass but no new candidates, and
random assembly produces short and repetitive walks far more often than long,
complete ones. What matters is the probability that the support contains a
witness, which Section~\ref{sec:scaling} computes.
\end{Pitfall}

Our executable model generates every legal walk containing between one and seven
vertices, starting from any vertex. The bound is explicit because the $B,C$
cycle allows arbitrarily long walks. There are 106 such bounded walks. This
finite enumeration is exhaustive within its declared domain; a prepared physical
population is a sample and does not inherit that guarantee.

We also inject one copy of $ABECDF$. This sequence is not generated by the ideal
edge-valid walk generator. It stands for a representation error that local
presence and length tests alone cannot exclude. Its role is logical fault
injection, not a quantitative model of contamination.

\section{Filtering: why the predicates fit together}\label{sec:filter}
Once the encoding is fixed, the intended logic is easy to state. For a candidate
$P=(v_1,\ldots,v_k)$ define
\[
 \begin{aligned}
 \mathcal E(P)&=[v_1=A\ \land\ v_k=F],\\
 \mathcal L(P)&=[k=6],\\
 \mathcal C(P)&=[\text{every vertex in }V\text{ occurs in }P].
 \end{aligned}
\]
Brackets denote a Boolean predicate, not a concentration. An ideal filter retains
exactly the candidates for which its predicate is true.

\begin{proposition}
If a sequence has exactly $n$ positions and contains every member of a set of $n$
vertices, each vertex occurs exactly once and no other symbol occurs.
\end{proposition}
\begin{proof}
Each of the $n$ required distinct vertices occupies at least one position. These
already use all $n$ positions. A repetition or an additional symbol would require
another position or displace a required vertex.
\end{proof}
The proposition explains why length and coverage together remove repeated
vertices. It says nothing about adjacency. A final edge check is still required
unless edge validity is guaranteed by the physical representation.

Figure~\ref{fig:pipeline} shows how each predicate is approximated by a
laboratory operation. Read each panel twice: once as chemistry (the picture) and
once as logic (the teal box underneath).

\TikzPlate{chapters/ch01/figures/workflow}{The molecular workflow of
generate--filter--verify, with each physical step paired with the predicate it
approximates. (1)~The graph is encoded as edge strands and vertex splints.
(2)~Annealing and ligation build random walks. (3)~PCR with primers at $A$ and
$F$ amplifies walks with the right endpoints. (4)~A gel keeps strands of the
length of six vertex codes. (5)~Beads carrying probes for $B$, $C$, $D$ and $E$
keep strands containing every vertex. (6)~Graduated PCR reveals the vertex order
and an independent verifier checks the edges. Dashed arrows are transfers of
material; each step has a finite yield and can retain wrong
molecules.}{fig:pipeline}

\begin{TryThis}
Before reading the table below, predict how many of the 106 legal walks survive
the endpoint filter, then the length filter, then the coverage filter. Which
walks survive coverage but would be rejected by an edge check? (Hint: none of
the legal walks; but the injected pseudo-path $ABECDF$ does.)
\end{TryThis}

The program applies the endpoint, length and coverage predicates, then
independently verifies the decoded route against $G$. Its exact output is:
\begin{center}\begin{tabular}{lrr}
\toprule Stage & Legal walks & With pseudo-path \\
\midrule
Generated & 106 & 107 \\
Endpoints & 14 & 15 \\
Length & 4 & 5 \\
Coverage & 2 & 3 \\
Verified & 2 & 2 \\
\bottomrule
\end{tabular}
\end{center}
Both populations finish with $ABCDEF$ and $ACBDEF$. The pseudo-path survives the
first three filters and is rejected only by full verification. The table makes a
useful distinction visible: \emph{passing selection is not the same event as
satisfying the specification}.

Physical selection is not Boolean. If a valid molecule survives stage $i$ with
probability $q_i$, conditional on surviving the preceding stages, its overall
survival probability is $\prod_i q_i$ by the chain rule of probability
(Appendix~\ref{app:math}). Independence is unnecessary for these conditional
probabilities. If one instead measures each stage's yield in isolation and
multiplies those numbers, additional assumptions are needed. Six hypothetical
stages each retaining 90 per cent leave $0.9^6\approx 53.1$ per cent overall.
This is a sensitivity calculation, not an estimate of Adleman's yield.

\section{Laboratory mechanisms and readable evidence}\label{sec:readout}
\Term{Hybridization} is association of complementary strands. \Term{Ligation}
makes a covalent backbone bond between suitable adjacent ends: T4 DNA ligase seals
a nick between a $3'$ hydroxyl and a $5'$ phosphate in a duplex.\cite{NEBT4Ligase}
Pairing and ligation therefore do different jobs: one aligns, the other connects.
A diagram that shows only base-pair contacts has not yet shown a continuous
product strand.

The \Term{polymerase chain reaction} (PCR) repeatedly separates strands and uses
primers and a polymerase to copy a selected region.\cite{NHGRIPCR} Primers
complementary to the endpoint codes make endpoint-compatible candidates
preferentially amplifiable. This is enrichment by chemistry, not a perfect
software predicate: amplification of an undesired template can also produce a
strong signal.

\Term{Gel electrophoresis} separates charged molecules in a gel under an electric
field; under suitable conditions fragment size largely determines migration.\cite{NHGRIElectrophoresis}
Selecting a size fraction implements a length constraint only as accurately as
the separation and recovery allow. In an affinity step, a complementary probe
recognises a required subsequence and retains molecules carrying it. Recognition
does not reveal the complete order of the strand. Table~\ref{tab:mechanisms}
summarises the correspondence and its characteristic failures.

\begin{table}[h]
\centering\small
\caption{Each ideal predicate is approximated by a physical mechanism with its own
failure modes. Measured specificity, yield and detection sensitivity are needed
before a physical step can be treated as its ideal predicate.}\label{tab:mechanisms}
\begin{tabular}{@{}p{0.2\linewidth}p{0.36\linewidth}p{0.36\linewidth}@{}}
\toprule
Predicate & Intended physical mechanism & Representative failure\\
\midrule
join edges & splint pairing, then ligase seals the nick & mismatched or unjoined ends\\
$\mathcal E$: endpoints & primers at $A$ and $F$, PCR & off-target amplification\\
$\mathcal L$: length & gel size selection & overlapping or smeared bands\\
$\mathcal C$: contains $v$ & capture by a complementary probe & loss of targets; nonspecific retention\\
witness & decoded order plus independent verifier & ambiguous signal from mixed molecules\\
\bottomrule
\end{tabular}
\end{table}

For readout, consider a primer anchored at the start code and a second primer
associated with an internal vertex. In our idealised six-vertex, 20-bases-per-vertex
model, the corresponding product lengths locate that vertex in a pure witness.
Figure~\ref{fig:readout} draws the expected lengths for $ABCDEF$.
\DeepFigure{DNAD-01-F8}{Synthetic graduated-PCR-style length readout for the
teaching witness. Lanes $B$ through $F$ show expected products of 40, 60, 80, 100
and 120 bases under the idealised encoding. These bands are neither photographed
data nor a reconstruction of the historical gel.}{fig:readout}

Now mix $ABCDEF$ with $ACBDEF$. Both $B$ and $C$ can produce lengths associated
with positions two and three. The union of bands records possibilities across
molecules; it does not assign every band to one molecule. Recovering a complete
witness requires resolving that ambiguity, for example by isolating single
molecules and reading them, followed by checking the decoded order.

\begin{Pitfall}
A band is a population average over many molecules, not a molecule. Two
different mixtures can produce the same band pattern (Chapter~10). A positive,
independently checked witness establishes existence; the \emph{absence} of a
detectable signal does not establish nonexistence without a model of generation,
retention and detection. Adleman's report itself discusses imperfect separation
and spurious paths as practical concerns.\cite{Adleman1994}
\end{Pitfall}

\section{Parallelism and computational complexity}\label{sec:complexity}
A candidate Hamiltonian route is a short \emph{certificate}. Given an adjacency
lookup and a table of vertices, a verifier scans the $n$ entries, rejects
duplicates or unknown vertices, checks the endpoints and checks $n-1$ edges. This
takes $O(n)$ lookups. Building an adjacency matrix may itself take $O(n^2)$ space;
representation and set-up belong in the resource account.

\begin{CSBridge}
\emph{Complexity classes in one paragraph.} P is the class of decision problems
that a standard computer can solve in time bounded by a polynomial in the input
size. NP is the class whose ``yes'' instances have a short certificate that can
be checked in polynomial time. A problem is NP-hard if every problem in NP can
be translated into it efficiently; it is NP-complete if it is also in NP.
Directed Hamiltonian path is NP-complete.\cite{PfenningHamilton} No one knows
whether P equals NP. These definitions do not say that every instance is hard.
Chapter~3 develops them carefully.
\end{CSBridge}

A tube can contain many candidates undergoing the same operation at once. This
changes the organisation of the work: material holds many possibilities
simultaneously. But a small number of filtering \emph{steps} says nothing about
the amount of \emph{material} needed to make the right candidate likely to exist.
Complexity accounting must include population size, representation length,
preparation, physical time and readout.

\section{Combinatorics becomes a material budget}\label{sec:scaling}
Suppose the endpoints are fixed in a complete directed graph. The $n-2$ interior
vertices can occur in
\[
 N=(n-2)!
\]
orders. This counts candidate orders; it does not say that a sparse graph
produces each one. Our nine-edge graph removes many orders through adjacency;
the complete-graph count isolates how quickly a naive permutation space grows.
\begin{center}\begin{tabular}{rr}
\toprule Vertices $n$ & Internal orders $(n-2)!$ \\
\midrule
6 & 24 \\
10 & 40,320 \\
15 & 6,227,020,800 \\
20 & 6,402,373,705,728,000 \\
25 & 25,852,016,738,884,976,640,000 \\
\bottomrule
\end{tabular}
\end{center}

Assume independent draws from some candidate-generation distribution, and let $p$
be the probability that one draw produces a chosen target, such as one
particular witness or any valid witness. After $M$ draws the probability of
missing it is $(1-p)^M$. Requiring a miss probability of at most $\delta$, where
$0<\delta<1$, gives
\[
 (1-p)^M\le\delta
 \quad\Longleftrightarrow\quad
 M\ln(1-p)\le\ln\delta
 \quad\Longleftrightarrow\quad
 M\ge\frac{\ln\delta}{\ln(1-p)}\quad(0<p<1).
\]
The denominator $\ln(1-p)$ is negative, which reverses the inequality when we
divide by it. The minimum integer $M$ is the ceiling of the right-hand side. When
$p=1$ one draw suffices; when $p=0$ no finite number does.

\DeepFigure{DNAD-01-F9}{Factorial counts of candidate orders on a logarithmic
vertical axis, together with the independent-sampling bound. Connecting segments
guide the eye; no measured chemical scaling law is fitted. Neither curve assigns
a concentration, reaction time or experimental yield.}{fig:scaling}

\begin{Worked}{how many molecules for 99 per cent confidence?}
For a uniform draw from the $4!=24$ interior orders of our six vertices, one
designated order has $p=1/24$. With $\delta=0.01$,
\[
M\ge\frac{\ln 0.01}{\ln(23/24)}=\frac{-4.6052}{-0.042560}=108.2,
\]
so 109 draws are needed. Our graph has two witnesses, so the event \emph{either
witness} has $p=2/24$ in the same hypothetical space and needs
$\lceil\ln0.01/\ln(11/12)\rceil=\lceil52.9\rceil=53$ draws. These are different
questions. The actual walk generator is neither uniform over permutations nor a
physical random sampler, so these numbers illustrate the method, not the
experiment.
\end{Worked}

For small $p$, $\ln(1-p)\approx-p$, giving $M\approx\ln(1/\delta)/p$: the number of
molecules needed grows in proportion to $1/p$, and $1/p$ grows factorially with
$n$. If each generated target survives the filters independently with
probability $q$, the effective per-draw success probability is $pq$. This
additional model must be declared rather than smuggled into a copy count.

Molecule count connects to a physical budget through
\[
 M=cV_{\ell}N_{\mathrm A},
\]
where $c$ is the amount concentration in moles per litre, $V_{\ell}$ is the volume
in litres and $N_{\mathrm A}=6.02214076\times10^{23}\,\mathrm{mol}^{-1}$ is the
exact Avogadro constant.\cite{BIPMmole} The units cancel to a count. Raising $M$
by increasing concentration or volume does not guarantee unchanged reaction
behaviour.

A useful resource description is therefore a vector, not a single speed number:
\[
 R=(\text{sequence inventory},\text{molecule count},\text{volume},
       \text{time},\text{energy},\text{readout cost}).
\]
Energy comparisons must specify whether synthesis, temperature control,
separation and measurement are included. Parallelism relocates costs; it does not
erase them. Chapter~4 develops this resource ledger in detail.

\section{The executable logical companion}\label{sec:code}
The implementation in \texttt{book/chapters/ch01/reference.py} deliberately avoids
an optimised solver. A bounded walk generator exposes repeated vertices, staged
filtering exposes which constraints do the work, and an independent permutation
oracle checks the final answer. The excerpt below is generated from the actual
source, so its line numbers and operations cannot silently diverge from the tested
program.
\VerbatimInput[fontsize=\footnotesize,numbers=left,numbersep=5pt,firstline=15,lastline=19]{chapters/ch01/reference.py}

\VerbatimInput[fontsize=\footnotesize,numbers=left,numbersep=5pt,firstline=21,lastline=30]{chapters/ch01/reference.py}

\VerbatimInput[fontsize=\footnotesize,numbers=left,numbersep=5pt,firstline=32,lastline=43]{chapters/ch01/reference.py}


There are two independent routes to the expected survivors. The stage model
starts with all bounded legal walks. The oracle fixes the endpoints, permutes the
four interior vertices and checks adjacency. Agreement on this graph catches
omissions in either construction more effectively than hard-coding one expected
route.

Other tests target boundaries: a repeated vertex can pass length, a pseudo-path
can pass coverage, an empty route cannot be a witness, and duplicated molecular
candidates should remain duplicated until an explicitly set-valued result is
requested. Converting a route to a set too early destroys precisely the order and
multiplicity the problem asks us to reason about.

The sampling helper checks its domain and is tested at the integer threshold: $M$
must meet the requested miss bound while $M-1$ must fail it. None of these tests
establishes hybridisation specificity or reaction yield. They establish the
consistency of an exact, inspectable logical model that a physical implementation
must then refine.

\section{From one search to a field}\label{sec:frontier}
DNA computing did not stop at generate-and-filter search. Benenson and colleagues
built an autonomous finite-state machine from DNA and enzymes.\cite{Benenson2001}
Soloveichik, Seelig and Winfree showed that DNA strand displacement can in
principle implement any chemical reaction network.\cite{Soloveichik2010} Qian and
Winfree built multilayer strand-displacement logic circuits, including one that
computes the square root of a four-bit number,\cite{QianWinfree2011} and Winfree and colleagues demonstrated designed
two-dimensional DNA crystals that self-assemble.\cite{Winfree1998} More recently,
one reprogrammable set of DNA tiles ran 21 different algorithms,\cite{Woods2019}
DNA neural networks learned to classify patterns from molecular
examples,\cite{CherryQian2025} and a scaffolded DNA computer whose answer is the
most thermodynamically favoured assembly ran programs such as multi-bit
addition.\cite{Sterin2026} Figure~\ref{fig:frontier} places these developments in
four overlapping eras.

\TikzPlate{chapters/ch01/figures/eras}{Four overlapping eras of DNA computing.
Search in a tube filtered molecular libraries; molecular machines and
self-assembly made computation autonomous; strand displacement let designers
program reaction dynamics directly in sequence; and recent work emphasises
reprogrammable hardware, molecular learning and integration with storage and
living cells. Each era brought its own representation, mechanism and
correctness argument; the dates mark when a style of work became prominent, not
sharp boundaries.}{fig:frontier}

The common idea is programmable physical organisation. The hard questions then
become specific. Can local recognition support reliable composition? Which
errors accumulate with depth or population size? How do we restore signals,
supply free energy, isolate unwanted interactions, and obtain an answer without
spending more effort on readout than the computation saves? The rest of this
book answers these questions one mechanism at a time.

Adleman's example gives us a disciplined starting vocabulary: specification,
encoding, population, transformation, selection, observation and verification.
Chapter~2 reconstructs the founding experiment closely and examines what its
controls and observations support. Part~II then supplies the chemistry and
physics deliberately abstracted here.

\begin{KeyIdeas}
\begin{itemize}\setlength{\itemsep}{2pt}
\item A DNA computation has three levels: a formal specification, an imperfect
physical process, and an interpretation of observations. Correctness is defined
at the first, achieved approximately at the second and argued at the third.
\item Strands have direction. Edge strands $R_uL_v$ and antiparallel vertex
splints turn graph adjacency into molecular adjacency; ligase turns adjacency
into one covalent strand.
\item A tube holds a \emph{multiset}: copy number, number of distinct candidates and
coverage of solutions are different quantities.
\item Length plus coverage implies no repetition, but not legal adjacency;
an independent verifier is part of the algorithm.
\item Each laboratory filter approximates a predicate with finite yield and
imperfect specificity. A detected witness proves existence; no signal does not
prove absence.
\item To include a target with probability $1-\delta$ needs about
$\ln(1/\delta)/p$ molecules. Because $p$ shrinks factorially, molecular
parallelism moves exponential cost into material; it does not remove it.
\end{itemize}
\end{KeyIdeas}

\section{Exercises}\label{sec:exercises}
\begin{enumerate}
\item Verify $ACBDEF$ against every condition in the definition. List the five
directed edges used.
\item Prove that length plus coverage excludes repetition. Give a six-symbol
counterexample showing that it does not enforce adjacency.
\item Draw the domain-level junction at $C$ for $B\to C\to D$. Mark both strand
directions. Compute the reverse complement of ACGA.
\item Identify the first ideal filter that removes $ABDEF$ and the first that
removes $ABCBDF$, assuming endpoint filtering comes first.
\item Add ten copies of $ABECDF$ to the 106 legal walks. Predict every stage
count without running the program; then verify.
\item Under uniform sampling of all fixed-endpoint interior orders, derive the
draw count for a 99 per cent chance of obtaining either witness.
\item Five stages each retain a target with conditional probability 0.8.
Calculate overall retention. Explain why isolated-stage measurements might not
justify the same product.
\item Design separate logical controls for candidate generation, selection loss
and readout failure. State what each can and cannot distinguish.
\item Compare an empty result from exhaustive symbolic enumeration with an empty
physical readout. Which additional assumptions would justify the same
conclusion?
\item Propose a fair resource comparison between a molecular implementation and
a digital solver. Specify the instance family, success criterion and accounting
boundary.
\end{enumerate}

\section{Worked solutions and discussion}\label{sec:solutions}
\begin{enumerate}
\item The edges are $AC,CB,BD,DE,EF$. All belong to $E$; the six vertices are
distinct, with endpoints $A,F$. Thus it is the second witness.
\item The $n$ required identities occupy all $n$ slots, leaving no room for a
duplicate. But $ABECDF$ has length six and complete coverage while $BE,EC\notin E$.
\item The upper fragments are $b_Rc_L$ and $c_Rd_L$, both written $5'\to3'$. A
complementary $c_L^{*}c_R^{*}$ splint runs $3'\to5'$ underneath and aligns the
nick between $c_L$ and $c_R$. ACGA has reverse complement TCGT.
\item Both pass endpoints. $ABDEF$ fails length. $ABCBDF$ passes length but fails
coverage because $E$ is absent.
\item Counts are $116,24,14,12,2$ at generation, endpoints, length, coverage and
verification. All ten injected copies persist until the edge check.
\item $p=2/24=1/12$; $\lceil\ln(0.01)/\ln(11/12)\rceil=53$. This is a
sampling-model calculation, not a prediction of molecular yield.
\item Retention is $0.8^5=0.32768$. The conditional formulation already accounts
for the selected population reaching each stage. Isolated yields may differ after
earlier processing changes that population.
\item A known distinctive target introduced before generation tests a different
chain from one introduced immediately before readout. Controls entering between
selection stages can help localise loss. Their behaviour need not match every
unknown candidate, so they bound an inference rather than prove universal
recovery.
\item Complete symbolic enumeration with a correct verifier establishes that no
witness exists in its finite search domain. A physical negative additionally
needs adequate coverage, recovery and detection sensitivity, with an explicit
failure bound.
\item Use identical mathematical instances and success requirements; include
preparation, failed runs, decoding and verification. Report distributions of
total cost and success, not only the fastest central operation. No winner follows
from the proposal alone.
\end{enumerate}

\section*{Further reading}
\addcontentsline{toc}{section}{Further reading}
Adleman's four-page report remains the best introduction to the founding
idea.\cite{Adleman1994} Lipton extended the approach to satisfiability,\cite{Lipton1995}
and Boneh, Dunworth, Lipton and Sgall analysed what such filtering models can
compute.\cite{Boneh1996} For the broader field, the review by Zhang and Seelig
introduces dynamic DNA nanotechnology,\cite{ZhangSeelig2011} and Seeman and
Sleiman survey DNA as an engineering material.\cite{SeemanSleiman2018}
